% =========================================================
% capitulos/07_Conclusiones.tex
% Placeholder con las herramientas que se van a usar en la
% tesis: tablas con el estilo azul, pseudocódigo y ecuaciones
% numeradas / sin numerar de distintos tipos.
% =========================================================

\chapter{Conclusiones}

\label{chap:Conclusiones}

Se logro desarrollar e implementar el marco de trabajo completo mismo que
constituye una alternativa híbrida viable para construir proyecciones discriminantes no
lineales mediante expansión polinomial explícita y selección evolutiva.

Su desempeño no muestra una superioridad sistemática frente a métodos kernell del
estado del arte en términos de exactitud, pero sí evidencian un comportamiento competitivo
en diversos escenarios y una ventaja metodológica relevante: la posibilidad de obtener
representaciones compactas, proyectables y trazables.

Los resultados sugieren que la elección del método de clasificación debe entenderse
como un compromiso entre la complejidad del problema, la capacidad predictiva requerida
y el nivel de interpretabilidad deseado. Mientras que en algunos escenarios los métodos
clásicos continúan siendo la alternativa más eficiente, en otros resulta necesario recurrir a
modelos capaces de representar relaciones no lineales más complejas. GNLPDA se posiciona
precisamente en ese punto intermedio, ofreciendo una solución que combina expansión
explícita de características, selección evolutiva y reducción de dimensionalidad con el
propósito de mantener un equilibrio entre rendimiento e interpretabilidad.









